
\section{Teoremas de comparación de los recubrimientos étales}
\label{XII.5}
\setcounter{subsection}{-1}
\subsection{}
\label{XII.5.0}
Precisemos la noción de recubrimiento finito de un espacio analítico.
\index{recubrimiento finito de un espacio analítico|hyperpage}%
Si $\othercal{X}$ es un espacio analítico, se dice que un espacio analítico
$\othercal{X}'$ finito sobre~$\othercal{X}$ es un recubrimiento finito de
$\othercal{X}$ si cada componente irreducible de~$\othercal{X}'$ domina
una componente irreducible de~$\othercal{X}$.


\begin{theoreme}[«Teorema de existencia de Riemann»]
\label{XII.5.1}
\index{Riemann (teorema de existencia de)|hyperpage}%
\index{teorema de existencia de Riemann|hyperpage}%
Sea $X$ un $\CC$-esquema localmente de tipo finito y sea $X^{\an}$
el espacio analítico asociado a~$X$. El funtor $\Psi$ que asigna
a cada recubrimiento étale finito $X'$ de~$X$ el espacio
$X^{\prime\an}$ es una equivalencia
\marginpar{333}
de la categoría de los recubrimientos étale finitos de~$X$ en
la categoría de los recubrimientos étale finitos de~$X^{\an}$.
\end{theoreme}

1) \emph{El funtor $\Psi$ es plenamente fiel.} Sean $X'$ y
$X''$ dos recubrimientos étale finitos de~$X$ y demostremos que
la aplicación canónica
\begin{equation*}
\label{eq:XII.5.*}
\tag{$*$} { \Hom_{X}(X',X'') \to
\Hom_{X^{\an}}(X^{\prime\an},X^{\prime\prime\an}) }
\end{equation*}
es biyectiva. Podemos suponer que $X'$ es conexo. Dar un
$X$-morfismo de~$X'$ en $X''$ equivale a dar una
componente conexa $X_i$ de~$X'\times_XX''$ tal que el morfismo
$X_i\to X'$ inducido por la primera proyección sea un
isomorfismo. Como las componentes conexas de~$X'\times_XX''$
corresponden biyectivamente a las componentes conexas de
$X^{\prime\an}\times_{X^{\an}}X^{\prime\prime\an}$~\eqref{XII.2.6}, y un
morfismo $X_i\to X'$ es un isomorfismo
\ifthenelse{\boolean{orig}}
{si y solo si}
{si y solo si}%
~lo es $X_i^{\an}\to X^{\prime\an}$, queda demostrada la
biyectividad de~\eqref{eq:XII.5.*}.

2) \emph{El funtor $\Psi$ es esencialmente sobreyectivo.} Sea
$\othercal{X}'$ un recubrimiento étale finito de~$X^{\an}$ y demostremos
que existe un recubrimiento étale $X'$ de~$X$ junto con
un isomorfismo $X^{\prime\an}\isomto\othercal{X}'$. En vista de~1), la
cuestión es local sobre~$X$ y podemos suponer, por tanto, que $X$ es afín.

a) Reducción al caso en que $X$ es normal. Podemos suponer que $X$
es reducido. Supongamos, en efecto, demostrado el teorema
para $X_{\red}$. El funtor que a un recubrimiento étale
finito $X'$ de~$X$ le asigna el recubrimiento étale finito
$X^{\prime\an}_{\red}$ de~$X_{\red}^{\an}$ es entonces una equivalencia.
Como se obtiene componiendo $\Psi$ con el funtor $\Theta$ que
a un recubrimiento étale finito de~$X^{\an}$ le asigna su imagen
inversa sobre~$X_{\red}^{\an}$, y $\Theta$ es plenamente
fiel, se sigue que $\Psi$ es una equivalencia de categorías.

Podemos
\marginpar{334}
suponer que $X$ es normal. Sea, en efecto, $\widetilde{X}$ la
normalización de~$X$ y sea $p\colon\widetilde{X}\to X$ el morfismo
canónico. Como $p$ es finito, $p$ es un morfismo de descenso
efectivo para la categoría de los recubrimientos
étales~(\Ref{IX.4.7}). Supuesto demostrado el teorema
para $\widetilde{X}$, pongamos
$\widetilde{\othercal{X}'}=\othercal{X}'\times_{X^{\an}}\widetilde{X}^{\an}$.
Existen un recubrimiento étale $\widetilde{X}'$ de
$\widetilde{X}$ y un isomorfismo $\widetilde{X}'{}^{\an}\simeq
\widetilde{\othercal{X}'}$. Resulta entonces de~1) que el dato
de descenso natural sobre~$\widetilde{\othercal{X}'}$ se levanta
a un dato de descenso sobre~$\widetilde{X}'$ relativo a
$\widetilde{X}\to X$; esto prueba la existencia de un recubrimiento
étale $X'$ de~$X$ y de un isomorfismo
$i\colon X^{\prime\an}\times_{X^{\an}}\widetilde{X}^{\an}\simeq
\widetilde{\othercal{X}'}$, cuyas imágenes inversas por las dos
proyecciones de
$\widetilde{X}^{\an}\times_{X^{\an}}\widetilde{X}^{\an}$ son iguales.
Según~\Ref{IX.3.2}, cuya demostración vale en el caso analítico,
el morfismo $\widetilde{X}^{\an}\to X^{\an}$ es un morfismo de descenso
para la categoría de los recubrimientos étales y, por consiguiente,
$i$ procede de un isomorfismo $X^{\prime\an}\simeq\othercal{X}'$.

b) Reducción al caso en que $X$ es regular. Sean $U$
el abierto de los puntos regulares de~$X$, $i\colon U\to X$ e
$i^{\an}\colon U^{\an}\to X^{\an}$ los morfismos canónicos; como
$X$ es normal, se tiene $\codim(X-U,X)\geq2$. Supongamos que existe un
recubrimiento étale $U'$ de~$U$ tal que
$U^{\prime\an}\simeq\othercal{X}'|U^{\an}$ y demostremos que entonces $U'$ se
prolonga en un recubrimiento étale $X'$ de~$X$ tal que
$X^{\prime\an}\simeq\othercal{X}'$. Basta ver que $U'$ se prolonga
en un recubrimiento étale $X'$ de~$X$; se tendrá entonces un
isomorfismo $X^{\prime\an}|U^{\an}\simeq\othercal{X}'|U^{\an}$. Pero si
$\cal{F}$ y $\cal{G}$ son los haces coherentes de álgebras
sobre~$\cal{O}_{X^{\an}}$ que definen, respectivamente, $\othercal{X}'$ y
$X^{\prime\an}$, la normalidad de~$X$ y la desigualdad
$\codim(X-U,X)\geq2$ implican que los morfismos canónicos
$$
\cal{F}\to i_*^{\an}(\cal{F}|U^{\an}) \qquad
\cal{G}\to i_*^{\an}(\cal{G}|U^{\an})
$$
son
\marginpar{335}
isomorfismos \cite[\no 3~prop\ptbl 4]{XII.11}. Se deduce que $\cal{F}$
y $\cal{G}$, y por tanto también $X^{\prime\an}$ y $\othercal{X}'$, son
isomorfos.

Sea $\varphi\colon X^{\an}\to X$ el morfismo canónico. Como el problema
de prolongar $U'$ a~$X$ es local sobre~$X$, basta demostrar que, para
cada punto $y$ de~$X^{\an}-U^{\an}$, el recubrimiento étale
$U'_{\varphi(y)}=U'\times_X\Spec\cal{O}_{X,\varphi(y)}$ de
$U_{\varphi(y)}=U\times_X\Spec\cal{O}_{X,\varphi(y)}$ se prolonga a
$\Spec\cal{O}_{X,\varphi(y)}$. Sea $H$ la
$\cal{O}_U$-álgebra coherente que define $U'$. El morfismo canónico
$$
\alpha\colon(i_*H)^{\an}\to i_*^{\an}(H^{\an})=\cal{F}
$$
define un morfismo de haces de módulos sobre
$\Spec\cal{O}_{X^{\an},y}$:
$$
\alpha_y\colon(i_*H)^{\an}_y\to\cal{F}_y\quoi,
$$
cuya restricción a
$U_y=U_{\varphi(y)}\times_{\Spec\cal{O}_{X,\varphi(y)}}
\Spec\cal{O}_{X^{\an},y}$ es un isomorfismo. Esto demuestra que
$H|U_y$ es trivial y, por tanto, que $U'_{\varphi(y)}$ se prolonga a
$\Spec\cal{O}_{X,\varphi(y)}$.

c) Caso en que $X$ es afín y regular. Sea
$$
j\colon X\to P
$$
una compactificación de~$X$, donde $P$ es un $\CC$-esquema
proyectivo y $j$ una inmersión abierta dominante. Gracias al
teorema de resolución de singularidades~\cite{XII.8}, se pueden
encontrar un esquema regular $R$ y un morfismo proyectivo
$r\colon R\to P$ tales que $r$ induzca un isomorfismo
$r^{-1}(X)\simeq X$ y $r^{-1}(X)$ sea el complementario en~$R$
de un divisor con cruzamientos normales. Sea
\marginpar{336}
$$
k\colon X\to R
$$
la inmersión canónica. Vamos a demostrar que existe un recubrimiento
finito normal~\eqref{XII.5.0} $\othercal{R}'$ de~$R^{\an}$ que prolonga
el recubrimiento étale $X^{\prime\an}$. Según la
Proposición~\Ref{XII.5.3} siguiente, tal recubrimiento es único; el
problema de prolongar $X^{\prime\an}$ es, por tanto, local sobre~$R^{\an}$
en un entorno de~$R^{\an}-X^{\an}$. Ahora bien, cada punto de
$R^{\an}-X^{\an}$ tiene un entorno abierto $\othercal{V}$ isomorfo a una
bola de un espacio afín $\othercal{E}^n$, tal que
$\othercal{V}-\othercal{V}\cap X^{\an}$ está definido por la anulación de
las primeras $p$ funciones coordenadas $z_1,\dots,z_p$, con
$0\leq p\leq n$. El grupo fundamental de
$\othercal{U}=\othercal{V}\cap X^{\an}$ es isomorfo a~$\ZZ^p$, y todo
recubrimiento étale de~$\othercal{U}$ es cociente de un recubrimiento de
la forma
$$
\othercal{U}''=\othercal{U}[T_1,\dots,T_p]/
(T_1^{n_1}-z_1,\dots,T_p^{n_p}-z_p)\quoi,
$$
donde los $n_i$ son enteros $>0$, por un subgrupo $H$ del grupo de
Galois $\ZZ/n_1\ZZ\times\cdots\times\ZZ/n_p\ZZ$ de~$\othercal{U}''$.
Ahora bien, $\othercal{U}''$ se prolonga en el recubrimiento regular
$$
\othercal{V}''=\othercal{V}[T_1,\dots,T_p]/
(T_1^{n_1}-z_1,\dots,T_p^{n_p}-z_p)\quoi,
$$
de~$\othercal{V}$ sobre el cual actúa $H$, y el cociente de
$\othercal{V}''$ por $H$ es la prolongación buscada.

La demostración concluye entonces gracias a~\Ref{XII.4.6}.
El recubrimiento $\othercal{R}'$ procede de un recubrimiento finito $R'$
de~$R$; la restricción de~$R'$ a~$X$ es un recubrimiento $X'$ de
$X$ tal que $X^{\prime\an}\simeq\othercal{X}'$, y según
\Ref{XII.3.1}(iii), $X'$ es un recubrimiento étale de~$X$.

\begin{corollaire}
\label{XII.5.2}
Sea
\marginpar{337}%
$X$ un $\CC$-esquema conexo localmente de tipo finito,
sea $\varphi\colon X^{\an}{\to}X$ el morfismo canónico y sea $x$ un
punto de~$X^{\an}$. Sea $\pi_1(X^{\an},x)$ el grupo fundamental del
espacio topológico $X^{\an}$ en~$x$, y sea $\pi_1(X,\varphi(x))$ el
grupo fundamental del esquema $X$ en~$\varphi(x)$
\textup{(\Ref{V.7})}. Entonces $\pi_1(X,\varphi(x))$ es canónicamente
isomorfo a la completación de~$\pi_1(X^{\an},x)$ para la topología de los
subgrupos de índice finito.
\end{corollaire}

En efecto, sea $\cal{C}$ la categoría de los recubrimientos étale finitos
de~$X^{\an}$, sea $F$ el funtor de~$\cal{C}$ en~$\Ens$ que a cada
recubrimiento étale finito $\othercal{X}'$ de~$X^{\an}$ le asigna el
conjunto de los puntos de~$\othercal{X}'$ situados sobre~$x$, y sea
$\widehat{\pi}_1(X^{\an},x)$ el grupo profinito asociado a~$\cal{C}$ y~$F$
como en~\Ref{V.4}. Como todo recubrimiento étale finito de~$X^{\an}$
es cociente del recubrimiento universal por un subgrupo de índice finito,
$\widehat{\pi}_1(X^{\an},x)$ no es sino la completación de
$\pi_1(X^{\an},x)$ para la topología de los subgrupos de índice finito.
El corolario resulta, pues, de~\Ref{XII.5.1} y~\Ref{V.6.10}.
